\section{Data}\label{sec:data}
\subsection{Sources and label rules}
All images come from official or author-published releases, and the rules were fixed before training.
Each image receives one of six mutually exclusive labels. NIH ChestX-ray14 \cite{wang2017chestxray8}
supplies Normal (``No Finding'', at most one image per patient), Pleural Effusion, Pulmonary Nodule/Mass
and Pneumonia (each only when no other finding is present). Shenzhen \cite{jaeger2014two} supplies Normal and
Tuberculosis. Kermany \cite{kermany2018identifying} supplies Normal and Pneumonia. The Cohen collection
\cite{cohen2020covid} supplies frontal COVID-19, Pneumonia, Normal and Tuberculosis images. TBX11K
\cite{liu2020rethinking} contributes healthy and active-TB images only; its ``sick but non-TB'' class (a
proxy that the earlier build had mislabelled as Pneumonia), latent TB and unlabelled images are excluded.
The COVID-19 Radiography Database \cite{chowdhury2020can,rahman2021exploring} contributes COVID-19 images
only. The final build (\RDatasetBuild) uses all six sources; the first V6 build (run 1) used the first four.
Lateral views and CT are excluded. Per-(source, class) caps bound CPU training time and limit any
source dominating a class.

\subsection{Patient-grouped split and leakage checks}
The final build has \RNTrain{} training, \RNVal{} validation and \RNTest{} test images, split 70/15/15 and
stratified by class and source. A group is a genuine or filename-derived patient identifier where one
exists (NIH, Kermany, Shenzhen, Cohen). For TBX11K and the COVID-19 Radiography Database patient identity is
unknown, so each image is its own group; the same patient could therefore appear in two splits for these
sources. Groups are merged through any shared exact (MD5) or near-duplicate (perceptual hash and thumbnail
correlation) image, so exact and near duplicates cannot cross splits. Separation of patients is therefore guaranteed only for
rows with known patient identity. Development images matching the external cohort are
excluded.

\subsection{External cohort}
The Montgomery County set \cite{jaeger2014two} (\RNExternal{} scans: \RExtSpecN{} normal, \RExtTBN{}
tuberculosis) were not used for training or for selecting the served model; however, run 2 was started after run 1's external results were seen (Section~\ref{sec:method}). The task is binary (Normal versus abnormal/TB), not a
six-class evaluation. Class and source remain partly confounded (COVID-19 images come from web
collections and Pneumonia partly from paediatric images), which matters for
Section~\ref{sec:results}.
